\section{Feuille TD1}\label{sec:td1}

% --- EXERCICE 1 ---
\begin{td-exo}[Appliquer \algo{AC}]
    On considère les réseaux de contraintes suivants :
    \begin{itemize}
        \item Soit \(N_1 = (X, D, C)\) avec \(X = \left\{X_1, X_2, X_3\right\}\) et 
        \begin{equation*}
            D = 
            \begin{cases}
                D(X_1) = \left\{1, 2, 3, 4, 5\right\} \\
                D(X_2) = \left\{1, 3, 5\right\} \\
                D(X_3) = \left\{1, 3, 4, 7\right\},
            \end{cases} \qquad
            C = 
            \begin{cases}
                X_1 + X_2  = X_3 \\
                X_1 < X_2.
            \end{cases}
        \end{equation*}

        \item Soit \(N_2 = (X, D, C)\) avec \(X = \left\{ X_1, \ldots, X_4\right\}\) et 
        \begin{equation*}
            \forall i \in \left\{1,\ldots 4\right\},
            D(X_i) = \left\{1,\ldots, 8\right\},
            \qquad
            C = 
            \begin{cases}
                2 X_1 = X_3 \\
                X_1 < X_2 \\
                X_3 + X_4 = 10 \\
                X_3 < X_2.
            \end{cases}
        \end{equation*}

        \item Soit \(N_3 = (X, D, C)\) avec \(X = \left\{ X_1, \ldots, X_5\right\}\) et 
        \begin{equation*}
            D = 
            \begin{cases}
                D(X_1) = \left\{1, 2, 7, 8\right\} \\
                D(X_2) = \left\{1, 2, 7, 8\right\} \\
                D(X_3) = \left\{1, 2, 3, 4\right\} \\
                D(X_4) = \left\{1, 2, 3, 4\right\} \\
                D(X_5) = \left\{0, 1\right\},
            \end{cases} \qquad
            C = 
            \begin{cases}
                X_1 + X_4 = 10 - 5X_5 \\
                \vert X_1 - X_2 \vert = \vert X_3 - X_4 \vert \\
                X_3 = 2X_2.
            \end{cases}
        \end{equation*}
    \end{itemize}

    \begin{enumerate}
        \item Appliquer \algo{AC} sur \(N_1\).
        \item Appliquer \algo{AC} sur \(N_2\).
        \item Appliquer \algo{AC} sur \(N_3\).
    \end{enumerate}
\end{td-exo}


\begin{td-sol} 
    \begin{enumerate}
        \item On applique \algo{AC} sur \(N_1\) : 
        \begin{equation*}
            \begin{aligned}
                c_2 &\implies (X_1, 5), (X_2, 1) \\
                c_1 &\implies (X_1, 3), (X_3, 1, 3).
            \end{aligned}
        \end{equation*}
        Après application de \algo{AC}, \(D\) devient :
        \begin{equation*}
            D = 
            \begin{cases}
                D(X_1) = \left\{1, 2, 4\right\} \\
                D(X_2) = \left\{3, 5\right\} \\
                D(X_3) = \left\{4, 7\right\}.
            \end{cases}
        \end{equation*}

        \item On applique \algo{AC} sur \(N_2\) :
        \begin{equation*}
            \begin{aligned}
                c_1 &\implies (X_3, 1, 3, 5, 7), (X_1, 5, 6, 7, 8) \\
                c_2 &\implies (X_2, 1) \\
                c_3 &\implies (X_4, 1, 3, 5, 7) \\
                c_4 &\implies (X_3, 8), (X_2, 2) \\
                c_3 &\implies (X_4, 2) \\
                c_1 &\implies (X_1, 4).
            \end{aligned}
        \end{equation*}
        Après application de \algo{AC}, \(D\) devient :
        \begin{equation*}
            D = 
            \begin{cases}
                D(X_1) = \left\{1, 2, 3\right\} \\
                D(X_2) = \left\{3, \ldots, 8\right\} \\
                D(X_3) = \left\{2, 4, 6\right\} \\
                D(X_4) = \left\{4, 6, 8\right\}.
            \end{cases}
        \end{equation*}

        \item On applique \algo{AC} sur \(N_3\) :
        \begin{equation*}
            \begin{aligned}
                c_1 &\implies (X_4, 1) \\
                c_3 &\implies (X_2, 7, 8), (X_3, 1, 3) \\
                c_2 &\implies (X_1, 7, 8) \\
                c_1 &\implies (X_5, 0), (X_4, 2).
            \end{aligned}
        \end{equation*}
        Après application de \algo{AC}, \(D\) devient :
        \begin{equation*}
            D = 
            \begin{cases}
                D(X_1) = \left\{1, 2\right\} \\
                D(X_2) = \left\{1, 2\right\} \\
                D(X_3) = \left\{2, 4\right\} \\
                D(X_4) = \left\{3, 4\right\} \\
                D(X_5) = \left\{1\right\}.
            \end{cases}
        \end{equation*}
    \end{enumerate}
\end{td-sol}

% --- EXERCICE 2 ---


\begin{td-exo}
    Soit le réseau de contraintes \(N = (X, D, C)\) où \(X = \left\{ X_1,\ldots X_n\right\}\) avec pour domaines :
    \begin{equation*}
        D = 
        \begin{cases}
            D(X_1) &= \left\{0, 1, 2, 3, 6\right\} \\
            D(X_2) &= \left\{1, 2, 3, 4\right\} \\
            D(X_3) &= \left\{0, 1, 3, 4, 5\right\} \\
            D(X_4) &= \left\{3, 4, 6, 10, 15\right\} \\
            D(X_5) &= \left\{0, 3, 4\right\}
        \end{cases}
    \end{equation*}
    et pour contraintes :
    \begin{equation*}
        C = 
        \begin{cases}
            X_1 + X_2 + X_3 = X_4 \\
            |X_1 - X_3| = X_5 \\
            |X_1 - X_2| = 2 \\
            X_3 + X_5 \neq X_1 + X_2 + X_4
        \end{cases}
    \end{equation*}
    Pour gagner du temps, on admettra que \(N\) est arc cohérent.
    \begin{enumerate}
        \item Appliquer l'algorithme MAC au réseau de contraintes \(N\).
        On utilisera l'heuristique dom/deg pour le choix des variables de l'ordre lexicographique pour le choix des valeurs.
        Si plusieurs variables ont le m\^eme ratio dom/deg, MAC choisira en priorité la variable \(X_i\) avec l'indice \(i\) le plus petit.
    \end{enumerate}
\end{td-exo}

\begin{td-exo}
    Soit le réseau de contraintes \(N = (X, D, C)\) où \(X = \left\{X_1,\ldots,X_4\right\}\) avec pour domaines :
    \begin{equation*}
        D = 
        \begin{cases}
            D(X_1) &= \left\{1, 2, 7, 8\right\} \\
            D(X_2) &= \left\{1, 2, 7, 8\right\} \\
            D(X_3) &= \left\{1, 2, 3, 4\right\} \\
            D(X_4) &= \left\{1, 2, 3, 4\right\}
        \end{cases}
    \end{equation*}
    et pour contraintes :
    \begin{equation*}
        C = 
        \begin{cases}
            X_1 + X_4 = 5 \\
            |X_1 - X_2| = |X_3 - X_4| \\
            X_3 = 2 X_2
        \end{cases}
    \end{equation*}
    \begin{enumerate}
        \item Appliquer BC sur \(N\).

        \item Appliquer SAC sur \(N\).

        \item Combien comporte de variables le plus petit réseau de contraintes binaires normalisé (c'est-à-dire tel que deux contraintes ne peuvent pas avoir la même portée) qui est SAC, mais n'a pas de solution ? Justifiez votre réponse.

        \item On peut définir la propriété SBC de façon analogue à SAC, en remplaçant AC par BC dans la définition. Comparez les propriétés AC, SAC, BC, SBC.\\
        Un réseau de contraintes binaires est \(k\)-quasi arborescent s'il est normalisé et qu'il existe \(k\) variables dont la suppression (ainsi que la suppression des contraintes incidentes) produit un réseau arborescent.

        \item Un réseau \(1\)-quasi arborescent qui est SAC admet-il toujours une solution ? Justifiez votre réponse.

        \item Soit \(k\geq 1\). Peut-on résoudre en temps polynomial les réseaux de contraintes \(k\)-quasi arborescents ? Justifiez votre réponse.
    \end{enumerate}
\end{td-exo}